\documentclass[runningheads]{llncs}

\usepackage[T1]{fontenc}
\usepackage{graphicx}
\usepackage{booktabs}
\usepackage{xcolor}
\usepackage{amsmath}

\newcommand{\todo}[1]{\textcolor{red}{[\textbf{TODO:} #1]}}
\newcommand{\tsc}[1]{\textsc{#1}}

\begin{document}

\title{Byzantine Cheap Talk: Adversarial Resilience and\\Topology Effects in LLM Coordination Games}

\author{Author~1\inst{1} \and Author~2\inst{1} \and Author~3\inst{1}}
\authorrunning{Author 1 et al.}
\institute{Affiliation \\ \email{placeholder@example.com}}

\maketitle

\begin{abstract}
Multi-agent LLM systems increasingly rely on communication protocols for coordination, yet their robustness under adversarial and structural constraints remains poorly understood.
Building on prior work showing that cheap-talk channels enable cooperation in LLM coordination games, we investigate two distinct vulnerability classes in a 4-player Stag Hunt.
First, Byzantine agents---who signal cooperation but defect---can eliminate group-level cooperation entirely, with honest agents detecting betrayal within one round but unable to recover coordination due to the game's all-or-nothing payoff structure; probabilistic adversaries partially preserve cooperation but extract the full cooperation surplus.
Second, explicitly restricting communication topology collapses cooperation, while identical restrictions applied silently preserve near-perfect cooperation---establishing that the mechanism is agents' meta-reasoning about hidden information, not information loss itself.
Finally, we document persistent behavioral archetypes across model families: models that defect rationally upon detecting betrayal vastly outperform models that persistently attempt cooperation under adversarial conditions, a gap that compresses under probabilistic deception.
Together, these findings have direct implications for the design of robust networked multi-agent AI systems.

\keywords{LLM agents \and game theory \and Byzantine fault tolerance \and cheap talk \and multi-agent coordination}
\end{abstract}

%===================================================================
\section{Introduction}\label{sec:intro}
%===================================================================

Large language models are increasingly deployed as autonomous agents in multi-agent systems where coordination is essential~\cite{wooldridge2009introduction,park2023generative}.
In classical game theory, \emph{cheap talk}---costless, non-binding pre-play communication---can resolve coordination dilemmas by establishing focal points~\cite{crawford1982strategic,farrell1987cheap}.
Recent work confirmed this for LLM agents: adding a single-word broadcast to a 4-player Stag Hunt raised cooperation from 0\% to 96.7\%~\cite{madmoun2025communication}.
However, this assumes honest signaling and unrestricted communication---conditions rarely met in practice.

We investigate two vulnerability classes.
First, \emph{Byzantine agents}~\cite{lamport1982byzantine}---who signal cooperation but defect---test whether honest agents can sustain coordination under strategic deception (Section~\ref{sec:byzantine}).
Second, \emph{communication topology restrictions} test whether partial message visibility preserves the coordination benefit (Section~\ref{sec:topology}).
Our experiments reveal that (1)~a single deterministic adversary eliminates group cooperation entirely, while probabilistic adversaries extract a 22\% payoff premium; (2)~topology restrictions collapse cooperation only when agents \emph{know} about them, not from information loss per se; and (3)~two stable behavioral archetypes---\emph{fast defectors} vs.\ \emph{persistent cooperators}---emerge across model families, with up to 16$\times$ payoff differences under adversarial conditions.

%===================================================================
\section{Background and Related Work}\label{sec:background}
%===================================================================

\paragraph{Cheap talk and coordination games.}
Crawford and Sobel~\cite{crawford1982strategic} showed that costless pre-play messages can transmit information even without binding commitments.
Farrell~\cite{farrell1987cheap} demonstrated that cheap talk selects among equilibria in coordination games.
The Stag Hunt~\cite{skyrms2004stag} is canonical: players prefer the payoff-dominant equilibrium (mutual cooperation) but face strategic risk, since unilateral cooperation yields the worst outcome.
Cheap talk can resolve this by signaling intent~\cite{farrell1996cheap}.

\paragraph{Byzantine fault tolerance.}
The Byzantine Generals Problem~\cite{lamport1982byzantine} formalized achieving consensus when some processes behave arbitrarily.
Lamport et al.~\cite{lamport1982byzantine} established that consensus tolerating $f$ Byzantine faults requires $n \geq 3f{+}1$ processes; practical protocols followed~\cite{castro1999practical}.
We adapt this framework to multi-agent AI: a Byzantine agent signals cooperation via cheap talk but defects at the action stage---strategic rather than arbitrary failure.

\paragraph{LLM multi-agent coordination.}
Game-theoretic evaluation of LLM agents is a growing area~\cite{sun2025survey}, with benchmarks revealing systematic behavioral patterns influenced by game framing~\cite{montrealethics_framing,huang2025gama_bench,guo2024investigation}.
Particularly relevant is the finding that contextual framing shifts LLM behavior between cooperative and competitive strategies~\cite{montrealethics_framing}---a theme recurring in our topology experiments.
We build directly on~\cite{madmoun2025communication}, who showed cheap talk raises Stag Hunt cooperation from 0\% to 96.7\% (corrected from initially reported 48.3\%).

\paragraph{Communication topology.}
Network structure fundamentally affects coordination in multi-agent systems~\cite{jackson2008social}.
We test whether LLM agents are robust to topology restrictions, and whether \emph{awareness} of restriction matters independently of the restriction itself.

%===================================================================
\section{Experimental Setup}\label{sec:setup}
%===================================================================

\paragraph{Game.}
We use a 4-player Stag Hunt with cheap talk~\cite{madmoun2025communication}.
Each of 5 rounds has two stages: (1)~agents simultaneously broadcast one word (non-binding cheap talk), then (2)~each chooses \tsc{Hunt~Stag} or \tsc{Hunt~Hare}.
The payoff for agent~$i$ is:
\begin{equation}\label{eq:payoff}
u_i = \begin{cases}
10 & \text{if all 4 choose \tsc{Stag},} \\
3 & \text{if } i \text{ chooses \tsc{Hare},} \\
0 & \text{if } i \text{ chooses \tsc{Stag} and any other chooses \tsc{Hare}.}
\end{cases}
\end{equation}
Agents observe the full history of prior rounds (all messages, actions, and payoffs).

\paragraph{Prompt structure.}
Each agent receives a system prompt specifying the game rules, its player identity, and the payoff structure.
At the communication stage, agents are instructed: ``\emph{Before you choose your action, you must broadcast ONE single word to the group. This word can be anything. It is non-binding cheap talk.}''
At the action stage, agents see all broadcast words and are asked to reason step-by-step before choosing, outputting structured JSON with a \texttt{reasoning} field and an \texttt{action} field.
\todo{Add GPT-4o prompt variant when available.}

\paragraph{Models.}
Four heterogeneous LLMs via DeepInfra API (temperature~0.7):
Mixtral-8x22B-Instruct~\cite{jiang2024mixtral},
Qwen2.5-72B-Instruct~\cite{qwen2024qwen25},
Llama-3.3-70B-Instruct~\cite{grattafiori2024llama},
DeepSeek-V3~\cite{deepseek2024v3}.
Fixed agent--model assignment across all trials.

\paragraph{Conditions.}
(i)~\textbf{Hard Byzantine} ($k \in \{0,1,2\}$): $k$ agents always broadcast ``stag'' but choose \tsc{Hare}; 10 trials each.
(ii)~\textbf{Soft Byzantine} ($k{=}1$, $p{=}0.5$): adversary defects with probability 0.5; 15 trials.
(iii)~\textbf{Explicit topology}: broadcast/ring/star with agents told their visibility constraints; 10 trials each.
(iv)~\textbf{Silent topology}: same filtering, no visibility cues, player count removed from prompt; 10 trials each.
(v)~\textbf{Byzantine $\times$ star}: star topology with 1 hard adversary, in two sub-conditions: adversary as hub vs.\ adversary as spoke; 10 trials each.

\paragraph{Metrics.}
\emph{Group cooperation}: fraction of rounds where all agents choose \tsc{Stag}. \emph{Honest cooperation}: fraction of honest \tsc{Stag} choices.

%===================================================================
\section{Byzantine Cheap Talk}\label{sec:byzantine}
%===================================================================

%-------------------------------------------------------------------
\subsection{Hard Byzantine Agents Eliminate Cooperation}\label{sec:byz-hard}
%-------------------------------------------------------------------

A hard Byzantine agent broadcasts ``stag'' but deterministically chooses \tsc{Hare}, making group cooperation structurally impossible (unanimity required).

\begin{table}[t]
\centering
\caption{Byzantine results across conditions (10 trials each, 5 rounds).}\label{tab:byzantine}
\smallskip
\begin{tabular}{@{}lccc@{}}
\toprule
Condition & Group\% & Honest\% & Avg.\ payoff \\
\midrule
Baseline ($k{=}0$)              & \todo{92.0} & \todo{95.0} & \todo{9.35} \\
Hard Byzantine ($k{=}1$)        & \todo{0.0}  & \todo{60.0} & \todo{1.20} \\
Hard Byzantine ($k{=}2$)        & \todo{0.0}  & \todo{37.0} & \todo{1.89} \\
Star + Byz.\ (hub{=}adversary)  & \todo{0.0}  & \todo{13.3} & \todo{2.60} \\
Star + Byz.\ (hub{=}honest)     & \todo{0.0}  & \todo{27.3} & \todo{2.18} \\
\bottomrule
\end{tabular}
\end{table}

While group cooperation drops to 0\% with a single adversary (Table~\ref{tab:byzantine}), 60\% of honest agents continue choosing \tsc{Stag} under $k{=}1$.
Chain-of-thought analysis confirms agents detect the message--action mismatch within one round, but cannot recover: the all-or-nothing payoff structure means even one defector renders cooperation futile.
Under $k{=}1$, honest cooperation starts at 100\% in round~1 (agents trust the initial cheap-talk signals), drops sharply to 53\% in round~2 after observing the first betrayal, then stabilizes near 50\% for the remaining rounds.
Under $k{=}2$, initial cooperation (95\%) collapses to 20\% by round~2 and remains there---a faster and more complete erosion, as agents face two simultaneous sources of deception.

\paragraph{Behavioral archetypes.}
Two model-family-stable patterns emerge under $k{=}1$.
\emph{Fast Defectors} (FD: Mixtral, DeepSeek) permanently switch to \tsc{Hare} after one betrayal---all 14 instances switched after round~1.
\emph{Persistent Cooperators} (PC: Qwen, Llama) continue choosing \tsc{Stag} despite exploitation---12/16 instances never switched.
\todo{RESULTS: Insert switching behavior breakdown by model and $k$.}

The payoff gap is stark: FD agents earn \todo{2.40}/round vs.\ PC agents' \todo{0.15}/round (16$\times$).
Every FD instance earns exactly 12 points (0 from the betrayed first round, then $3{\times}4$ from hare-hunting).
Fast defection is individually rational but collectively destructive.

%-------------------------------------------------------------------
\subsection{Soft Byzantine: Probabilistic Deception}\label{sec:byz-soft}
%-------------------------------------------------------------------

A soft adversary ($p{=}0.5$) creates intermittent rather than systematic betrayal.
Group cooperation partially survives at \todo{21.3\%} (vs.\ 0\% hard), because the adversary sometimes cooperates.
Honest cooperation reaches \todo{71.1\%}, decaying gradually from \todo{97.8\%} in round~1 to \todo{51.1\%} in round~5 but never collapsing.
The adversary earns \todo{3.65}/round vs.\ honest agents' \todo{3.00}/round---a 22\% premium, extracting the cooperation surplus.

Critically, the FD/PC payoff gap compresses from 16$\times$ (hard) to 1.9$\times$ (soft).
When betrayal is intermittent, persistent cooperators occasionally achieve group success (earning 10 rather than 0), while fast defectors forfeit these windfalls by committing to hare-hunting.
The adversary's advantage is structural: by always signaling ``stag,'' it free-rides on rounds when it cooperates (sharing the 10-point stag payoff) while capturing the 3-point hare floor when it defects---honest agents cannot distinguish genuine cooperation from a lucky coin flip.

%-------------------------------------------------------------------
\subsection{Topology $\times$ Adversary Interaction}\label{sec:byz-star}
%-------------------------------------------------------------------

To test whether Byzantine deception and topology restriction interact, we place one hard adversary in a star topology under two conditions: the adversary is assigned as the hub (visible to all spokes) or as a spoke (visible only to the hub).
Both conditions yield 0\% group cooperation, but honest cooperation diverges sharply: \todo{13.3\%} when the adversary is the hub vs.\ \todo{27.3\%} when honest (Table~\ref{tab:byzantine}).
Hub position amplifies adversarial influence because the hub's ``stag'' signal reaches all spokes, and spokes---seeing only the hub---cannot cross-validate against other players' messages.
Cooperation collapses almost instantly when the adversary is hub (46.7\% in round~1, 0\% by round~2), while an honest hub sustains higher initial trust (73.3\% in round~1, decaying to 16.7\% by round~5).

Compared to broadcast Byzantine ($k{=}1$), where honest cooperation is 60\%, the star topology compounds adversarial damage to 13--27\%---the two vulnerability classes interact multiplicatively rather than independently.
The FD/PC payoff gap compresses further to 1.4--1.7$\times$, as star topology suppresses cooperation so heavily that archetype differences are muted.
\todo{Validate with GPT-4o and Claude Sonnet.}

%===================================================================
\section{Communication Topology}\label{sec:topology}
%===================================================================

%-------------------------------------------------------------------
\subsection{Explicit vs.\ Silent Topology}\label{sec:topo-results}
%-------------------------------------------------------------------

We test three topologies---broadcast (all-to-all), ring (two neighbors), star (center sees all; periphery sees center only)---under two conditions: \emph{explicit} (agents told their visibility constraints) and \emph{silent} (same filtering, no cues, player count removed).

\begin{table}[t]
\centering
\caption{Topology results (10 trials per condition, 5 rounds each). Silent conditions use leak-fixed v2 prompts.}\label{tab:topology}
\smallskip
\begin{tabular}{@{}lcccccc@{}}
\toprule
& \multicolumn{3}{c}{Explicit} & \multicolumn{3}{c}{Silent} \\
\cmidrule(lr){2-4} \cmidrule(lr){5-7}
& Broadcast & Ring & Star & Broadcast & Ring & Star \\
\midrule
Group coop.   & \todo{80\%}  & \todo{0\%}  & \todo{0\%}  & \todo{100\%} & \todo{100\%} & \todo{100\%} \\
Agent coop.   & \todo{87.5\%} & \todo{33.5\%} & \todo{27\%} & \todo{100\%} & \todo{100\%} & \todo{100\%} \\
Avg.\ payoff  & \todo{8.38}  & \todo{2.00} & \todo{2.19} & \todo{10.0} & \todo{10.0} & \todo{10.0} \\
\bottomrule
\end{tabular}
\end{table}

Explicit ring and star collapse cooperation to 0\% group rate (Table~\ref{tab:topology}), matching hard Byzantine severity despite no adversary being present.
The FD/PC archetype split reappears in explicit topologies: in the ring, Llama cooperates at 72\% and Qwen at 56\%, while Mixtral (2\%) and DeepSeek (4\%) defect almost universally---the same model-family alignment as under Byzantine conditions.

The silent conditions produce a strikingly different outcome: all three topologies achieve perfect cooperation (100\% group, 100\% individual) with zero defections across 600 total agent-round observations.
The FD/PC split vanishes entirely---all four model families cooperate at 100\% across all three topologies.
Note that in the silent ring, each agent sees only two of three co-players' messages, yet cooperates as reliably as under broadcast.
This rules out information loss as the driver of the explicit topology collapse.

%-------------------------------------------------------------------
\subsection{Meta-Reasoning as the Mechanism}\label{sec:topo-mechanism}
%-------------------------------------------------------------------

The explicit--silent contrast establishes that cooperation collapse is not caused by information loss---agents receiving fewer messages cooperate perfectly when unaware of the restriction.
The mechanism is \emph{meta-reasoning}: told ``you can only see messages from your neighbors,'' agents infer unseen players with unknown intentions, triggering strategic uncertainty.
Without such framing, agents take messages at face value.
This parallels the framing effects documented by Lor\`e and Heydari~\cite{montrealethics_framing}: the same game-theoretic situation produces opposite outcomes depending on how information is presented.
It also reveals that the FD/PC archetypes are not intrinsic model properties but are \emph{elicited by strategic framing}---they emerge only when agents are prompted to reason about incomplete information, and disappear entirely when that prompt cue is absent.

%===================================================================
\section{Discussion and Conclusion}\label{sec:discussion}
%===================================================================

Cheap-talk coordination among LLM agents, while effective under benign conditions~\cite{madmoun2025communication}, is fragile along two axes.
\emph{Byzantine vulnerability}: a single deterministic adversary eliminates cooperation entirely---more severe than classical BFT, where $n{\geq}3f{+}1$ suffices~\cite{lamport1982byzantine}, because the Stag Hunt's unanimity requirement makes any defection catastrophic.
Probabilistic adversaries are more insidious: they sustain enough cooperation to extract a persistent \todo{22\%} surplus.
\emph{Topology vulnerability}: cooperation collapses only when agents \emph{know} communication is restricted---a cognitive failure (meta-reasoning about hidden information), not an informational one.
This implies that multi-agent systems exposing network topology may inadvertently undermine cooperation.

The FD/PC behavioral split is stable across all adversarial conditions yet vanishes entirely under silent topology, confirming these archetypes are prompt-elicited rather than reflecting intrinsic model dispositions.
No current model achieves the ideal response: conditional cooperation that detects deception and adapts without either abandoning cooperation prematurely (fast defectors) or persisting irrationally (persistent cooperators).
Designing agents---or prompts---that achieve this middle ground remains an open challenge.

\paragraph{Design implications.}
These findings suggest that robust multi-agent LLM systems require mechanisms beyond cheap talk: commitment devices, reputation systems, or verified communication channels.
Equally important, system designers should attend to \emph{how} agents are informed about their communication environment---exposing topology details can be as destructive as an active adversary.

\paragraph{Limitations and future work.}
All experiments use a single game (4-player Stag Hunt) with 10--15 trials.
The unanimity payoff structure may amplify adversarial impact; graded-payoff games may differ.
\todo{Extend to GPT-4o, Claude Sonnet, and IPGG+P.}

\bibliographystyle{splncs04}
\bibliography{refs}

\end{document}
